\section{Scope, antecedents, and the exact meaning of the results}
\label{scope:section}

The starting point is the canonical ProveIt manuscript
\cite{RepoMain} and the eleven archives present in \src{docs/incoming}
at repository commit
\begin{center}
\texttt{5a790187c8e186e41e2b990b4941cb7a1a3c7b6b}.
\end{center}
The commit timestamp is 11 October 2026, 00:09:09 UTC. The incoming
reports retain their 10 October dates. The delivered source manifest
records the retrieved paths and hashes. This article is an integration
candidate; it does not modify the canonical manuscript or promote
unreviewed incoming statements by editorial fiat.

The immediate motivation is three explicit unfinished parts of the
latest exact-identity programme. The ordered-Hurwitz report
\cite{RepoOrdered} asks for the first independent-slope formula beyond
depth three and for the directions where its additional ordering
coefficients disappear. The triple-Stieltjes report \cite{RepoTriple}
provides finite closure for reflected factors with distinct seams, while
coincidences require a separate normalization. The Gaussian depth
projector report \cite{gf:Depth} proposes completing the Cayley ideal
before repeating a finite reduction search. Sections~\ref{oq:section},
\ref{ct:section}, and~\ref{gf:sec} answer these questions in the
precise senses stated below. The harmonic-power work in
Section~\ref{hp:sec:main} gives a further exact family connecting
multiple polylogarithms, Gamma coefficients, harmonic numbers, and
spectral finite parts.

\subsection{A ledger of contributions}

\begin{longtable}{@{}p{0.22\textwidth}p{0.39\textwidth}p{0.32\textwidth}@{}}
\toprule
Direction & Proved here & Boundary of the claim\\
\midrule
\endhead
Depth-four ordered resonance & A spanning normal form for all transverse
straight rays, convergent coordinates, shifts, and reconstruction &
No arithmetic independence; no prescription for a ray contained in a
polar hyperplane.\\[5pt]
Coincident reflected products & All Stieltjes and argument derivative
indices; arbitrary finite reflected products; explicit Gamma primitives &
The generic unreflected cubic Tornheim coordinate is not evaluated.\\[5pt]
Raw harmonic powers & All nonpositive principal and finite coefficients
as shift polynomials; centered cancellation; explicit families &
The first regular spectral coefficient generally needs global Mellin
data.\\[5pt]
Gaussian formal obstruction & Nonmembership after adding the complete
Cayley ideal to the frozen additional relation span &
$S_6$ and revised $S_8$ remain conjectural as numerical period identities.\\[5pt]
Literature correction & A sign correction to equation (22) of
arXiv:2507.03058v1, with an all-logarithmic proof &
It is a correction to that identified formula, not a rejection of the
preprint's entire framework.\\
\bottomrule
\end{longtable}

The word ``new'' in this report means an explicitly demonstrated
continuation beyond the inspected ProveIt results. We identify the
classical inputs and the nearest primary-source antecedents. We do not
claim an exhaustive priority search over every formulation in the
literature. In particular, the Gamma diagonal of ordered regular germs,
ordinary Hurwitz Fourier expansions, shuffle regularization, and
raw-power meromorphic continuation are antecedents.

The analytic context is the Laurent theory of multiple zeta functions
developed by Saha \cite{Saha} and by Matsumoto, Onozuka, and Wakabayashi
\cite{MOW}. Hurwitz integral methods and their Gamma applications have
substantial classical development, including Espinosa and Moll
\cite{EspinosaMoll1,EspinosaMoll2}. Tin \cite{hp:Tin2025} studies raw
harmonic-power Dirichlet series and their harmonic Stieltjes constants;
Karg\i n, Dil, Cenkci, and Can \cite{Kargin} study related harmonic zeta
constants. We retain their strict/inclusive and inner/outer shift
distinctions rather than silently identifying different series.
The shuffle-polynomial and Eulerian-logarithm structure in the formal
argument is classical \cite{gf:Radford,gf:Patras}; the new certificate
concerns its application to the complete ideal intersection.

\subsection{Three normalizations that must not be interchanged}

For $a>0$, throughout the article,
\[
 \zeta(1+u,a)=\frac1u+
       \sum_{m\ge0}\frac{(-1)^m}{m!}\gamma_m(a)u^m,
 \qquad \gamma_0(a)=-\psi(a).
\]
A spectral derivative differentiates the order $s$ of a zeta function;
an argument derivative differentiates its positive shift $a$ or its
position $x$. A superscript $(p)$ on $\gamma_m$ or $\psi$ in the
reflected section always denotes an argument derivative. Bracketed
superscripts $[j]$ on $z_k$ in the ordered section denote spectral
derivatives. The ordinary constants are $\gamma_m=\gamma_m(1)$.
Generalized harmonic numbers are
$H_n^{(r)}=\sum_{j=1}^n j^{-r}$, with $H_0^{(r)}=0$.

There are three different finite-coefficient operations.
\begin{enumerate}
\item A \emph{spectral finite part} is the constant Laurent coefficient
in a specified complex parameter. The ray variable $\varepsilon$ in
Section~\ref{oq:section} and the spectral variable $s+m$ in
Section~\ref{hp:sec:main} are explicitly fixed.
\item A \emph{coordinate finite-part integral} is the constant term of
a cut-off integral in fixed local coordinates. On $(0,1)$ those
coordinates are $x$ and $1-x$. Powers and logarithms are expanded before
the constant term is selected. This is the convention in
Section~\ref{ct:section}.
\item A \emph{formal shuffle relation} is an equality in a specified
rational word algebra. Its truth and exact verification are independent
of an unproved assertion that the period map is injective.
\end{enumerate}
All ordinary integrals, convergent sums, analytic continuations, and
finite parts are identified where they first occur. Analytic continuation
of an integral and coefficientwise coordinate completion agree away
from a resonance, but can differ by a finite Taylor correction at one.
That correction is computed explicitly in the coincident theorem.

The multiple-polylogarithm convention used for the Gaussian target is
\[
 \operatorname{Li}_{r,s}(u,v)
 =\sum_{n>m\ge1}\frac{u^nv^m}{n^rm^s}.
\]
The one-variable multiple polylogarithms in the harmonic-power kernel
carry the argument only on the outer index. These conventions fix the
word order, colors, and harmonic conversions; changing them changes
intermediate coordinate formulas.

\subsection{Reading and integration map}

The four mathematical sections can be read independently after the
normalizations above. The ordered section supplies all of its required
lower-depth germs. The reflected section proves the completed transform
before using polynomial reduction. The harmonic section gives both a
finite general algorithm and closed low-weight examples. The formal
section includes the algebraic argument connecting its finite exact
checks to the unrestricted ideal. Section~\ref{audit:section} distinguishes
an actual correction from normalization safeguards, and
Section~\ref{research:section} identifies concrete remaining targets.

For canonical integration, the ordered section naturally follows the
latest ordered-Hurwitz contribution; the reflected section follows the
coincident/contact calculus; the harmonic section belongs with harmonic
polylogarithm and endpoint regularization material; and the complete
Cayley certificate belongs immediately after the existing bounded
$S_6$ obstruction. The retained proof-status qualifications are part of
each mathematical statement.
